\chapter{Preambuła}

Luka kompetencyjna na styku ocen rozwiązań informatycznych przez zespoły prawne oraz techniczne jest przyczyną nieprawidłowej klasyfikacji systemów informatycznych. 
Co więcej, organizacje nieprawidłowo zarządzają łańcuchem dostawców i poddostawców, nawet gdy są do tego zobowiązane przez rozporządzenie DORA oraz dyrektywę NIS2. Szczególnie jest to widoczne w relacjach z hiperskalerami objętymi jurysdykcją Cloud Act, gdzie bardzo często pomija się ryzyka transferu danych do USA, co stoi w sprzeczności z RODO.

Ocena Ryzyka Informatycznego Podmiotów Nadzorowanych, w skrócie zwana ORION, jest nową metodą pozwalającą na wykrywanie i adresowanie wskazanych luk przez organizacje.
Metoda adresowana jest przede wszystkim do podmiotów objętych nadzorem regulacyjnym, instytucji finansowych, operatorów usług kluczowych oraz dostawców usług cyfrowych, dla których zgodność z wymogami DORA, NIS2 i RODO stanowi obowiązek prawny, a nie dobrowolny standard. Warto jednak wykorzystać medodę ORION również w pozostałych przedsiębiorstwach, traktując podnoszone przez metodologię wyzwania wynikające z implementacji DORA oraz NIS2 jako zestaw dobrych praktych, równocześnie upewniając się, że działania podmiotu są zgodne z obowiązującymi przepisami np. w zakresie RODO.

Niniejsza publikacja przedstawia założenia teoretyczne metody ORION, opisuje jej strukturę klasyfikacyjną oraz ilustruje jej zastosowanie praktyczne na przykładzie wybranej organizacji.

\chapter{Wstęp}

Transformacja cyfrowa przedsiębiorstw i instytucji publicznych, postępująca intensywnie od końca XX wiek,  zmieniła charakter ryzyk i wyzwań stawianych przed organizacjami. Systemy informatyczne przestały być wyłącznie narzędziami wspierającymi codzienną pracę, stały się jej elementem strategii przedsiębiorstw, bez którego realizacja procesów jest praktycznie niemożliwa. Dotyczy to w równym stopniu banku przetwarzającego miliony transakcji konsumentów, szpitala prowadzącego elektroniczną dokumentację medyczną pacjentów, operatora infrastruktury energetycznej zarządzającego siecią dystrybucyjną, uczelni i szkoły prowadzącej cyfrowe ewidencje postępów podopiecznych, firm produkcyjnych opierających ciągłość działania na cyfrowych algorytmach. We wszystkich tych przypadkach niedostępność systemu informatycznego awarią o ograniczonych skutkach, jest zdarzeniem operacyjnym, prawnym często niosącym negatywne skutki finansowe.

Ustawodawca zdefiniował standardy bezpieczeństwa opisane w przepisach prawa. Do kluczowych aktów prawnych, które wymagają prowadzenia szczegółowej analizy ryzyka, należą przede wszystkim Ustawa o Krajowym Systemie Cyberbezpieczeństwa \autocite{uksc2018} oraz Rozporządzenie w sprawie operacyjnej odporności cyfrowej sektora finansowego, potocznie zwane DORA \autocite{dora2022}. Choć powyższe dokumenty jasno określają, kogo dotyczą ich zapisy, warto potraktować te publikacje jako uniwersalne źródło wiedzy. Wskazane w nich procedury oraz opisane wektory zagrożeń mogą pomóc zabezpieczyć każdą organizację przed skutkami cyberataków.

Równolegle z rosnącym uzależnieniem organizacji od technologii informatycznych rozwinął się złożony rynek dostawców, poddostawców i partnerów technologicznych, tworzący łańcuchy zależności ICT (\textit{Information and Communications Technology}). Pojedyncza aplikacja biznesowa może dziś opierać się na wielu elementach dostarczanych przez szereg dostawców: platformach chmurowych, bibliotekach oprogramowania, usługach uwierzytelniania, systemach płatności, zewnętrznych centrach przetwarzania danych czy dostawcach łączności. Każde z tych elementów jest potencjalnym źródłem ryzyka, technicznego (awarii), prawnego (brak zgodności z obowiązującym prawem) i operacyjnego. To organizacja ponosi odpowiedzialność regulacyjną za skutki zdarzeń w całym tym łańcuchu, niezależnie od tego, czy posiada nad nim bezpośrednią kontrolę. Złożoność łańcuchów dostaw ICT sprawia, że ich pełna inwentaryzacja, ocena i nadzór są zadaniem, z którym większość organizacji radzi sobie niewystarczająco,  nie z braku dobrej woli, lecz z braku odpowiednich narzędzi metodycznych.

Problem ten nabiera szczególnego wymiaru w kontekście intensywnej aktywności legislacyjnej europejskiego prawodawcy ostatniej dekady. Ochrona danych stanowi jeden z fundamentalnych elementów współczesnego porządku prawnego Unii Europejskiej. Dynamiczny rozwój technologii, w szczególności sztucznej inteligencji oraz gromadzenia  informacji, a co za tym idzie szeroka cyfryzacja usług publicznych i prywatnych oraz rosnąca skala międzynarodowego przetwarzania danych osobowych sprawiły, że zagadnienie ochrony danych stało się przedmiotem prac legislacyjnych zarówno na poziomie unijnym, jak i krajowym. Unia Europejska oraz jej państwa członkowskie zwracają szczególną uwagę na jakość ochrony danych, uznając ją za istotny element ochrony praw jednostki, zaufania do gospodarki oraz bezpiecznego funkcjonowania społeczeństwa informacyjnego. Regulacje te, omówione szczegółowo w rozdziale \textit{Krajobraz regulacyjny}, tworzą środowisko, w którym zgodność z prawem stała się warunkiem koniecznym prowadzenia działalności w sektorach takich jak finanse, ochrona zdrowia, energetyka czy transport. Niespełnienie wymogów wiąże się z sankcjami finansowymi wobec przedsiębiorstw, a odpowiedzialność za naruszenia spoczywa bezpośrednio na organach zarządczych podmiotów nadzorowanych.

W tym kontekście szczególnego znaczenia nabiera pytanie badawcze niniejszej monografii: w jaki sposób organizacja może w sposób wiarygodny, powtarzalny i oparty na obiektywnych dowodach ocenić rzeczywisty stan zgodności swoich systemów informatycznych i ich łańcucha dostawców z obowiązującymi wymogami prawnymi? Jest to teoretycznie proste pytanie. Praktyka audytowa i badawcza pokazuje bowiem, że istniejące metodyki, choć wartościowe w swoich dziedzinach, nie dają na nie zadowalającej odpowiedzi. Powstała więc luka, której wymiary nie były do tej pory szacowane, jednak luka jest realna, dotychczas nieadresowana, jednocześnie wskazuje na braki prawidłowej oceny zgodności zarządzania systemami informatycznymi.
Deklarowana zgodność w wymagami regulatorów rynku często rozmija się z weryfikowalnym stanem faktycznym. Organizacja posiadają certyfikaty i polityki bezpieczeństwa, lecz nie potrafią wskazać, które konkretne systemy przetwarzają dane szczególnej kategorii, kto jest faktycznym podddostawcą kluczowej platformy chmurowej ani czy dostawca zewnętrzny kiedykolwiek był przedmiotem oceny zgodności. Zjawisko to, opisywane w literaturze anglojęzycznej jako \textit{compliance theater}, wynika ze słabości systemów oceny opartych wyłącznie na deklaracjach i samoocenie. Jest ono równie szkodliwe, jak jawna niezgodność: daje organizacji fałszywe poczucie bezpieczeństwa.

\section{Postulat: Ocena Ryzyka Nadzorowanych Systemów Informatycznych }
Prezentowana praca opisuje metodologię Oceny Ryzyka Informatycznego Organizacji Nadzorowanych (ORION), oryginalny framework metodyczny zaprojektowany w celu wypełnienia opisanej luki. ORION nie jest kolejną implementacją istniejących standardów, nie zastępuje audytów ISO czy NIST, nie jest też systemem informatycznym w technicznym sensie tego słowa. Jest formalnie zdefiniowaną metodyką oceny zgodności systemów ICT i ich łańcucha dostaw, opartą na trzech innowacjach konceptualnych.

Fundamentem metodywki jest wieloperspektywiczność, którą można wytłumaczyć jako segmentacje kompetencyjną. Celem jest rozdzielenie procesu oceny na płaszczyzny: prawną, biznesową (użytkową) oraz techniczną. Zapewnia to, że odpowiedzi udzielają wyłącznie osoby posiadające realne kompetencje w danej dziedzinie.

Kolejny istotny element to rygor dowodowy. Kluczem jest zastąpienie deklaratywności ocen weryfikacją artefaktów (logi, certyfikaty). System wprowadza mechanizm penalizacji scoringowej za brak materialnych dowodów, co eliminuje zjawisko fasadowej zgodności (compliance theater).

Finalnie wynik mapowany do jako sieć powiązań w modelu grafowym. Dzięki algorytmom propagacji ryzyka, ORION pozwala wyliczyć, jak błędy poddostawcy wpływają na bezpieczeństwo i ciągłość działania podmiotu nadrzędnego.
